\chapter{Preface}

\section{Purpose of This Script}

This script is designed as a learning-oriented companion to the course
\emph{Theoretical Physics 2: Quantum Mechanics}.  It is not meant to replace the
lecture notes, exercise sheets, or official course material.  Its purpose is to
organize the material into a coherent sequence of concepts, derivations,
examples, and exam-oriented problem patterns.

The guiding principle is simple: every formula should be connected to a physical
interpretation, every formal statement should be connected to a calculational
method, and every standard calculation should be connected to typical exam
questions.

\begin{conceptbox}[title=How to use this script]
Read each chapter in three passes.
\begin{enumerate}[label=\arabic*.]
    \item First, read the definitions and conceptual explanations to understand
    what the objects mean.
    \item Second, work through the derivations and examples with pen and paper.
    \item Third, use the exam checklist at the end of the chapter to test whether
    you can reproduce the central arguments without looking them up.
\end{enumerate}
\end{conceptbox}

\section{Main Sources}

The main source for this script is the official lecture material for
\emph{Theoretical Physics 2: Quantum Theory}.  The handwritten lecture notes are
used as a supplementary source whenever they clarify the order of presentation
or contain useful intermediate calculations.  The problem sheets and their
solutions are used primarily to identify exam-relevant techniques and recurring
calculation patterns.

The source priority used in this script is:
\begin{enumerate}[label=\arabic*.]
    \item official exam information, exam sheets, topic lists, and learning goals;
    \item official lecture slides or lecture script;
    \item exercise sheets and official sample solutions;
    \item personal submitted solutions, only when they are readable and physically
    consistent;
    \item supplementary textbooks.
\end{enumerate}

\begin{notebox}[title=Source-critical use of solutions]
Personal solutions and sample solutions are not copied blindly.  They are used as
orientation for notation, level of detail, and typical problem-solving strategies.
If a derivation appears incomplete, ambiguous, or inconsistent, the script gives a
clean derivation instead.
\end{notebox}

\section{Recommended Literature}

The following books are listed in or are naturally aligned with the course
material.  They are useful in different ways:

\begin{itemize}
    \item F. Schwabl, \emph{Quantenmechanik (QM I): Eine Einfuehrung}.  Useful as a
    German-language reference with a structure close to many theoretical physics
    lectures.
    \item C. Cohen-Tannoudji, B. Diu, and F. Laloe, \emph{Quantum Mechanics},
    Volumes 1 and 2.  Very detailed and especially useful for conceptual
    explanations, examples, and complementary material.
    \item G. Grawert, \emph{Quantenmechanik}.  A compact German-language study text.
    \item J. J. Sakurai, \emph{Modern Quantum Mechanics}.  Particularly useful for
    the abstract Hilbert-space formalism, Dirac notation, spin, and angular
    momentum.
    \item A. Galindo and P. Pascual, \emph{Quantum Mechanics I and II}.  A more
    advanced and mathematically careful reference.
\end{itemize}

\section{Structure of the Script}

The script is organized so that the formalism is built gradually.  The early
chapters introduce wave functions, Hilbert spaces, probability interpretation,
and the Schrödinger equation.  The middle chapters develop operators,
commutators, measurements, uncertainty relations, and standard one-dimensional
systems.  Later chapters cover angular momentum, spin, perturbation theory,
identical particles, and elementary quantum information concepts such as qubits,
Bloch vectors, and entanglement.

\begin{center}
\begin{tabular}{@{}clp{8.5cm}@{}}
\toprule
No. & File & Topic \\
\midrule
00 & \texttt{00\_preface.tex} & Purpose, sources, learning strategy \\
01 & \texttt{01\_wave\_functions\_and\_statistical\_interpretation.tex} & Wave functions, Hilbert space, probability interpretation \\
02 & \texttt{02\_schroedinger\_equation\_and\_time\_evolution.tex} & Schrödinger equation, time evolution, probability current \\
03 & \texttt{03\_free\_particles\_wave\_packets\_and\_fourier\_representation.tex} & Free particles, plane waves, Fourier representation, spreading \\
04 & \texttt{04\_one\_dimensional\_potentials\_and\_tunneling.tex} & One-dimensional potentials, barriers, tunneling, scattering \\
05 & \texttt{05\_delta\_potential\_bound\_states\_and\_discontinuities.tex} & Delta potential, bound states, jump conditions \\
06 & \texttt{06\_operators\_commutators\_and\_ehrenfest\_theorem.tex} & Operators, commutators, Ehrenfest theorem \\
07 & \texttt{07\_measurements\_observables\_and\_uncertainty.tex} & Measurements, observables, compatible measurements, uncertainty \\
08 & \texttt{08\_dirac\_notation\_and\_abstract\_hilbert\_space.tex} & Dirac notation, projectors, matrix elements, representations \\
09 & \texttt{09\_harmonic\_oscillator\_and\_ladder\_operators.tex} & Harmonic oscillator and ladder operators \\
10 & \texttt{10\_coherent\_states\_and\_quasiclassical\_motion.tex} & Coherent states and quasicalssical motion \\
11 & \texttt{11\_angular\_momentum\_and\_spherical\_harmonics.tex} & Orbital angular momentum and spherical harmonics \\
12 & \texttt{12\_central\_potentials\_and\_hydrogen\_atom.tex} & Central potentials and the hydrogen atom \\
13 & \texttt{13\_scattering\_theory\_partial\_waves\_and\_resonances.tex} & Scattering theory, partial waves, resonances \\
14 & \texttt{14\_spin\_half\_systems\_and\_pauli\_matrices.tex} & Spin-1/2 systems and Pauli matrices \\
15 & \texttt{15\_time\_dependent\_spin\_dynamics\_and\_rabi\_formula.tex} & Rotating frame, magnetic resonance, Rabi formula \\
16 & \texttt{16\_addition\_of\_angular\_momenta\_and\_spin\_coupling.tex} & Addition of angular momenta and spin coupling \\
17 & \texttt{17\_time\_independent\_perturbation\_theory.tex} & Non-degenerate and degenerate perturbation theory \\
18 & \texttt{18\_identical\_particles\_bosons\_fermions\_and\_symmetrization.tex} & Identical particles, bosons, fermions, symmetrization \\
19 & \texttt{19\_qubits\_bloch\_sphere\_and\_entanglement.tex} & Qubits, Bloch sphere, Bell states, entanglement \\
20 & \texttt{20\_exam\_strategy\_formula\_sheet\_and\_problem\_patterns.tex} & Exam strategy, formula sheet, problem patterns \\
\bottomrule
\end{tabular}
\end{center}

\section{Notation and Conventions}

Unless stated otherwise, states are denoted either as wave functions
\(\psi(x)\), \(\psi(\vec{x})\), or as abstract vectors \(\ket{\psi}\).  Operators
are denoted by hats when helpful, for example \(\op{H}\), \(\op{X}\), and
\(\op{P}\).  The inner product may be written either as an integral in position
space or in Dirac notation:
\begin{equation}
    \braket{\phi}{\psi}
    = \int \dd x\, \phi^*(x)\psi(x).
\end{equation}

For a single particle in one spatial dimension, the canonical position and
momentum operators are usually represented by
\begin{equation}
    (\op{X}\psi)(x) = x\psi(x),
    \qquad
    (\op{P}\psi)(x) = -\ii\hbar \frac{\dd}{\dd x}\psi(x),
\end{equation}
and they satisfy the canonical commutation relation
\begin{equation}
    \comm{\op{X}}{\op{P}} = \ii\hbar\id.
\end{equation}

\begin{examtipbox}[title=Notation matters]
Many mistakes in quantum mechanics are notation mistakes.  Always distinguish
between a state, a wave function, an operator, an eigenvalue, and a measurement
result.  For example, \(\psi(x)\) is a representation of a state, while
\(\ket{\psi}\) is the abstract state itself.
\end{examtipbox}

\section{How to Study Problem Sheets}

The problem sheets are not merely exercises after the theory.  They define the
calculational core of the course.  For each problem, the following questions are
useful:
\begin{enumerate}[label=\arabic*.]
    \item Which postulate or principle is being tested?
    \item Which representation is most convenient: position space, momentum
    space, matrix notation, or Dirac notation?
    \item Which boundary conditions or normalization conditions are required?
    \item Which quantities are physically observable?
    \item Which limiting cases should be checked?
\end{enumerate}

\begin{conceptbox}[title=Typical quantum-mechanical workflow]
A large fraction of problems follows the same pattern:
\begin{enumerate}[label=\arabic*.]
    \item Identify the Hilbert space and the relevant operator.
    \item Solve the eigenvalue problem or expand the state in a useful basis.
    \item Normalize states and impose boundary conditions.
    \item Compute probabilities, expectation values, or transition amplitudes.
    \item Interpret the result physically and check limiting cases.
\end{enumerate}
\end{conceptbox}

\section{Common Sources of Errors}

\begin{itemize}
    \item Confusing \(\psi\) with \(\abs{\psi}^2\).  The wave function is a complex
    amplitude; only the modulus squared is a probability density.
    \item Forgetting normalization factors in continuous bases, especially in
    Fourier transforms.
    \item Treating non-normalizable plane waves as physical states without using
    wave packets or generalized normalization.
    \item Applying boundary conditions only to \(\psi\), but forgetting the
    derivative condition when the potential is finite.
    \item Applying derivative continuity at a delta potential, where the
    derivative has a jump instead.
    \item Ignoring degeneracy when using perturbation theory.
    \item Mixing up active time evolution of states with passive changes of basis.
\end{itemize}

\section{Summary}

This chapter sets the purpose and structure of the script.  The script is based
primarily on the official lecture material and the exercise sheets, with
handwritten notes and personal solutions used only as supplementary material.
The goal is to produce a coherent learning document that connects concepts,
derivations, examples, and exam-style problem solving.

\section{Exam Checklist}

Before moving on, make sure that you can answer the following questions:
\begin{itemize}
    \item What is the purpose of this script, and how should it be used?
    \item Which sources have priority when lecture notes, exercises, and personal
    solutions differ?
    \item What is the difference between a state, a wave function, an operator,
    and an observable?
    \item Why are normalization, boundary conditions, and basis choice central in
    quantum-mechanical calculations?
    \item Which chapters are likely to be especially relevant for exam-style
    calculations?
\end{itemize}
